\section{Coordinate Systems}
\label{sec:coordinate-systems}

A physical point or vector exists independently of the labels used to describe it. A coordinate system is therefore not part of nature itself; it is a mathematical convention that assigns numbers to geometric objects. Good coordinates make a problem transparent, while poor coordinates can hide its symmetry.

\subsection{Cartesian Coordinates}

In three-dimensional Euclidean space, a Cartesian coordinate system consists of three mutually perpendicular axes with associated orthonormal unit vectors
\[
\mathbf e_x,\qquad \mathbf e_y,\qquad \mathbf e_z.
\]
They satisfy
\[
\mathbf e_i\cdot\mathbf e_j=\delta_{ij},
\]

\begin{bookfigure}{A point and its Cartesian position vector. The point is geometric; the coordinates are numerical labels relative to the selected axes.}
\begin{tikzpicture}[scale=0.95]
  \coordinate (O) at (0,0);
  \coordinate (P) at (4.5,2.8);
  \draw[physics axis] (-0.25,0) -- (5.5,0) node[right] {$x$};
  \draw[physics axis] (0,-0.25) -- (0,3.8) node[above] {$y$};
  \draw[construction line] (P) -- (4.5,0);
  \draw[construction line] (P) -- (0,2.8);
  \draw[physics vector] (O) -- (P) node[midway,above left] {$\mathbf r$};
  \fill[visualgold] (P) circle (2.2pt);
  \node[above right] at (P) {$P(x,y)$};
  \node[below left] at (O) {$O$};
  \node[below] at (4.5,0) {$x$};
  \node[left] at (0,2.8) {$y$};
\end{tikzpicture}

\end{bookfigure}
where $\delta_{ij}$ is the Kronecker delta. A point $P$ is represented by the ordered triple $(x,y,z)$, and its position vector is
\[
\mathbf r=x\mathbf e_x+y\mathbf e_y+z\mathbf e_z.
\]

\begin{bookfigure}{The displacement from $P$ to $Q$ is $\Delta\mathbf r=\mathbf r_Q-\mathbf r_P$. Its Cartesian components form a right triangle.}
\begin{tikzpicture}[x=1cm,y=1cm]
\BookAxesTwoD{5.7}{3.9}
\coordinate (P) at (1.0,0.8); \coordinate (Q) at (4.7,3.0);
\draw[book guide] (P)--(4.7,0.8)--(Q);
\draw[book vector] (P)--(Q) node[midway,above left] {$\Delta\mathbf r=\mathbf r_Q-\mathbf r_P$};
\BookPoint{P}{below left:$P(x_1,y_1)$}
\BookPoint[fill=visualred]{Q}{above right:$Q(x_2,y_2)$}
\node at (2.85,0.53) {$\Delta x$}; \node at (4.98,1.9) {$\Delta y$};
\end{tikzpicture}

\label{fig:distance-between-points}
\end{bookfigure}

\begin{physical}
The point $P$ is geometric; the numbers $(x,y,z)$ are coordinate-dependent. Rotating the axes changes the components but does not move the point.
\end{physical}

\subsection{Coordinates and Components}

For a general vector $\mathbf A$,
\[
\mathbf A=A_x\mathbf e_x+A_y\mathbf e_y+A_z\mathbf e_z.
\]
The coefficients are obtained by projection:
\[
A_x=\mathbf A\cdot\mathbf e_x,\qquad
A_y=\mathbf A\cdot\mathbf e_y,\qquad
A_z=\mathbf A\cdot\mathbf e_z.
\]
Thus a basis converts an abstract vector into an ordered list of numbers.

\begin{bookfigure}{Rotating the coordinate axes changes the numerical components, but not the geometric vector itself.}
\begin{tikzpicture}[x=1cm,y=1cm]
\draw[book axis] (-.3,0)--(5.4,0) node[right] {$x$};
\draw[book axis] (0,-.3)--(0,4.1) node[above] {$y$};
\begin{scope}[rotate=28]
\draw[book axis,draw=visualcyan] (-.2,0)--(5,0) node[right] {$x'$};
\draw[book axis,draw=visualcyan] (0,-.2)--(0,3.5) node[above] {$y'$};
\end{scope}
\draw[book vector] (0,0)--(4.05,3.15) node[above right] {$\mathbf v$};
\draw[book angle] (1,0) arc[start angle=0,end angle=28,radius=1];
\node[text=visualgold] at (1.23,.27) {$\phi$};
\node at (3.9,.55) {same vector, new components};
\end{tikzpicture}

\label{fig:coordinate-rotation}
\end{bookfigure}

\booktopic{Other Useful Coordinate Systems}

Cartesian coordinates are not always optimal. Cylindrical coordinates $(\rho,\phi,z)$ are adapted to axial symmetry, while spherical coordinates $(r,\theta,\phi)$ are adapted to central symmetry. Their relations to Cartesian coordinates are
\[
x=\rho\cos\phi,\qquad y=\rho\sin\phi,
\]
and
\[
x=r\sin\theta\cos\phi,\quad
y=r\sin\theta\sin\phi,\quad
z=r\cos\theta.
\]
Unlike the Cartesian basis, curvilinear basis vectors generally vary from point to point.

\begin{examplebox}[Choosing coordinates]
The gravitational field of an isolated spherical body is most naturally expressed in spherical coordinates because the field depends primarily on the radial distance $r$. A long straight wire suggests cylindrical coordinates because rotations around the wire leave the physical situation unchanged.
\end{examplebox}

\begin{warningbox}
Coordinates are not vectors. The triple $(x,y,z)$ lists components relative to a chosen basis; the geometric position vector is $\mathbf r=x\mathbf e_x+y\mathbf e_y+z\mathbf e_z$.
\end{warningbox}
